%% GENERATED by tools/gen_latex.py -- edit content/agents09.py instead.
%% Build:  latexmk -pdf agents09.tex     (or: pdflatex agents09.tex, twice)
\documentclass[aspectratio=169,11pt,t]{beamer}
\usepackage{itbpro}

\coursetitle{Masukan yang bukan teks}
\coursesubtitle{Gambar, layar, dan suara --- apa yang berubah pada biaya, pada apa yang bisa jadi bukti, dan pada kewajiban yang menyertainya.}
\coursekicker{Chapter 9}
\coursesource{Grootendorst \& Alammar, An Illustrated Guide to AI Agents --- bab 9}
\courseduration{3 jam (2 sesi)}
\coursepresenter{Hendri Karisma}{Instructor}
\coursebrand{AI for Professional \textperiodcentered\ ITB \textperiodcentered\ Ch 9}

\title{Masukan yang bukan teks}
\author{Hendri Karisma}
\date{}

\begin{document}
\covertitle

\framekicker{Learning outcomes}
\begin{frame}[shrink=25]{Session Objectives}
\leadpar{By the end of this session, participants can:}
\begin{isteps}
\item \textbf{Menghitung biaya token sebuah gambar} dan menjelaskan kenapa ia tumbuh menurut luas.
\item \textbf{Menyebutkan tiga cara modalitas disambungkan} dan pertukarannya.
\item \textbf{Merancang alur dokumen} yang menalar atas data terstruktur, bukan atas piksel --- dan menyebutkan alasannya.
\item \textbf{Menyebutkan apa yang bisa dan tidak bisa diandalkan} dari model penglihatan pada dokumen nyata.
\item \textbf{Menjelaskan kenapa foto dokumen adalah data pribadi}, dan apa yang berubah pada rancangan karenanya.
\item \textbf{Menyebutkan bentuk bukti} yang membuat jawaban dari gambar bisa diperiksa manusia.
\end{isteps}

\vskip 6pt
\vskip 4pt
\reslink{Site}{Course home}{https://hendrikarisma.my.id/rs/ai-products-course/index.html}
\reslink{Book}{Maarten Grootendorst \& Jay Alammar, An Illustrated Guide to AI Agents (O'Reilly Media, Early Release)}{https://www.oreilly.com/library/view/an-illustrated-guide/9798341662681/}
\vskip 2pt
\end{frame}

\coursesection{01}{Apa yang sebenarnya berubah}{Gelungnya sama. Yang berubah biayanya, buktinya, dan kewajibannya.}

\framekicker{Dasar}
\begin{frame}[shrink=25]{Gelungnya tidak berubah sama sekali}
{\fontsize{9.4}{12.6}\selectfont\color{fg2}Agen yang menerima gambar tetap melakukan hal yang sama: memilih langkah, memanggil alat, membaca hasil, berhenti. Tidak ada mekanisme baru di Bab 1 sampai 8 yang jadi tidak berlaku.\par}\vskip 4pt
\begin{minipage}[t]{0.3253\textwidth}
\begin{icardaccent}
\cardh{Biayanya berubah}
Satu gambar bisa berharga ribuan token, dan ia masuk riwayat --- jadi dibayar ulang tiap giliran.
\end{icardaccent}
\end{minipage}%
\hfill
\begin{minipage}[t]{0.3253\textwidth}
\begin{icardaccent}
\cardh{Buktinya berubah}
\textbackslash\{\}u201cSaya membaca itu di gambar\textbackslash\{\}u201d tidak bisa diperiksa seperti hasil alat yang punya pengenal.
\end{icardaccent}
\end{minipage}%
\hfill
\begin{minipage}[t]{0.3253\textwidth}
\begin{icardaccent}
\cardh{Kewajibannya berubah}
Foto dokumen hampir selalu data pribadi, dan sering data yang diatur khusus.
\end{icardaccent}
\end{minipage}%
{\fontsize{9.4}{12.6}\selectfont\color{fg2}Ketiganya keputusan rekayasa, dan ketiganya diambil sebelum satu baris kode ditulis. Kedalaman arsitektur penglihatannya ada di \textbf{Bab 8--12} kelas ini.\par}\vskip 4pt
\end{frame}

\framekicker{Dasar}
\begin{frame}[shrink=25]{Cara sebuah gambar sampai ke model bahasa}
\begin{center}
\includegraphics[width=0.92\linewidth,height=0.60\textheight,keepaspectratio]{figs/agents09-encode.tex.pdf}
\end{center}
\vskip 3pt{\color{ink3}\fontsize{7.4}{9.4}\selectfont Gambar jadi vektor, disamakan ruangnya dengan teks, lalu diperlakukan sebagai token.\par}
{\fontsize{9.4}{12.6}\selectfont\color{fg2}Gambar dipecah jadi petak, tiap petak jadi vektor lewat penyandi penglihatan, lalu sebuah \textbf{penyambung} memetakan vektor itu ke ruang yang sama dengan token teks. Sesudah titik itu, model bahasa memperlakukan keduanya sama.\par}\vskip 4pt
\begin{iband}[signal]
Akibat praktis yang paling penting dari kalimat terakhir: \textbf{gambar menempati ruang jendela konteks yang sama dengan teks}, dan bersaing dengannya.
\end{iband}
\end{frame}

\coursesection{02}{Biayanya}{Satu gambar bukan satu lampiran.}

\framekicker{Biaya}
\begin{frame}[shrink=25]{Berapa token satu gambar}
\begin{center}
\includegraphics[width=0.97\linewidth,height=0.56\textheight,keepaspectratio]{figs/agents09-imgcost.tex.pdf}
\end{center}
\vskip 3pt{\color{ink3}\fontsize{7.4}{9.4}\selectfont Dihitung dari asumsi ubin yang tercetak di gambarnya. Langkahi menurut ukuran gambar.\par}
{\fontsize{9.4}{12.6}\selectfont\color{fg2}Satu foto dokumen berharga sekitar \textbf{3,3$\times$ halaman teks yang sama isinya} --- dan lebih sulit dibaca modelnya. Itu dua kerugian sekaligus, dan keduanya bisa dihindari kalau teksnya memang tersedia.\par}\vskip 4pt
\end{frame}

\framekicker{Biaya}
\begin{frame}[shrink=25]{Gambar masuk riwayat, dan riwayat dibayar ulang}
{\fontsize{9.4}{12.6}\selectfont\color{fg2}Bab 2 sudah menghitungnya untuk teks: tiap giliran membayar ulang semua giliran sebelumnya. Untuk gambar, angka yang dibayar ulang itu ribuan token, bukan ratusan.\par}\vskip 4pt
\begin{isteps}
\item \{'h': 'Keluarkan gambar dari riwayat begitu diekstrak', 'p': 'Simpan hasil ekstraksinya --- beberapa puluh token --- dan buang pikselnya dari konteks.'\}
\item \{'h': 'Turunkan resolusi ke yang benar-benar perlu', 'p': 'Empat kali token untuk resolusi dua kali lipat. Sebagian besar tugas tidak butuh resolusi penuh.'\}
\item \{'h': 'Potong ke bagian yang relevan', 'p': 'Satu kotak tanda tangan tidak butuh seluruh halaman.'\}
\end{isteps}
\begin{iband}[signal]
Langkah pertama sendirian sering memotong biaya alur dokumen \textbf{berkali lipat}, dan hampir tidak pernah dilakukan sebab gambar terasa seperti lampiran, bukan seperti token.
\end{iband}
\end{frame}

\coursesection{03}{Cara modalitas disambungkan}{Tiga pendekatan, dan pertukaran yang terasa di pemakaian.}

\framekicker{Penyambung}
\begin{frame}[shrink=25]{Tiga cara, dari yang paling sederhana}
\begin{center}
\includegraphics[width=0.92\linewidth,height=0.60\textheight,keepaspectratio]{figs/agents09-connect.tex.pdf}
\end{center}
\vskip 3pt{\color{ink3}\fontsize{7.4}{9.4}\selectfont Proyeksi, berbasis kueri, dan peleburan.\par}
{\fontsize{9.4}{12.6}\selectfont\color{fg2}Ketiganya menyelesaikan masalah yang sama --- membuat vektor gambar bisa dibaca model bahasa --- dan berbeda pada \textbf{berapa token yang dihasilkan} dan \textbf{apa yang hilang}.\par}\vskip 4pt
\end{frame}

\framekicker{Penyambung}
\begin{frame}[shrink=25]{\ldots{}dan pertukaran yang terasa di pemakaian}
\begin{itable}{>{\raggedright\arraybackslash}p{0.2256\textwidth}>{\raggedright\arraybackslash}p{0.2256\textwidth}>{\raggedright\arraybackslash}p{0.4888\textwidth}}
\thead{Cara} & \thead{Token gambar} & \thead{Pertukarannya} \\ \midrule
Proyeksi & Banyak, sebanding luas & Paling sederhana dan paling banyak dipakai; detail terjaga, biaya tinggi \\
Berbasis kueri & \textbf{Tetap} & Biaya bisa diramalkan; detail halus bisa terlewat karena jumlah tokennya tidak ikut naik \\
Peleburan & Bervariasi & Paling kuat pada tugas yang menuntut penalaran gabungan; paling mahal dilatih \\
\end{itable}
{\fontsize{9.4}{12.6}\selectfont\color{fg2}Untuk pembangun agen, baris kedua yang paling terasa: \textbf{biaya yang bisa diramalkan} kadang lebih berharga daripada ketepatan sedikit lebih tinggi.\par}\vskip 4pt
\end{frame}

\framekicker{Penyambung}
\begin{frame}[shrink=25]{Suara dan video: bentuk masalah yang sama}
\begin{minipage}[t]{0.485\textwidth}
{\fontsize{9.4}{12.6}\selectfont\color{fg2}\textbf{Suara}\par}\vskip 4pt
{\fontsize{9.4}{12.6}\selectfont\color{fg2}Disandikan jadi vektor lalu disambungkan dengan cara yang sama. Untuk hampir semua agen bisnis, jalur yang benar tetap: \textbf{transkripsikan dengan alat}, lalu bekerja atas teksnya.\par}\vskip 4pt

\end{minipage}%
\hfill
\begin{minipage}[t]{0.485\textwidth}
{\fontsize{9.4}{12.6}\selectfont\color{fg2}\textbf{Video}\par}\vskip 4pt
{\fontsize{9.4}{12.6}\selectfont\color{fg2}Bingkai demi bingkai, jadi biayanya berlipat dengan durasi. Sampel beberapa bingkai kunci hampir selalu cukup, dan selisih biayanya besar.\par}\vskip 4pt

\end{minipage}%
{\fontsize{9.4}{12.6}\selectfont\color{fg2}Pola yang sama muncul pada ketiga modalitas: \textbf{ubah jadi representasi yang lebih kecil dan lebih bisa diperiksa dengan alat, lalu nalarkan atas representasi itu.}\par}\vskip 4pt
\begin{iband}[signal]
Dan satu keuntungan yang jarang disebut: transkrip dan data terstruktur \textbf{bisa disimpan, dicari, dan dikutip}. Piksel dan gelombang tidak.
\end{iband}
\end{frame}

\framekicker{Penyandi}
\begin{frame}[shrink=25]{Kenapa penyandi gambar dan teks bisa bertemu sama sekali}
{\fontsize{9.4}{12.6}\selectfont\color{fg2}Pertanyaan yang wajar: bagaimana vektor dari piksel bisa berarti apa pun bagi model yang dilatih atas kata?\par}\vskip 4pt
{\fontsize{9.4}{12.6}\selectfont\color{fg2}Jawabannya pelatihan berpasangan. Kalau banyak sekali pasangan gambar dan keterangannya dilatihkan bersama, dengan tujuan \textbf{mendekatkan pasangan yang cocok dan menjauhkan yang tidak}, kedua penyandi lama-lama menaruh ``seekor kucing'' dan gambar kucing di tempat yang berdekatan.\par}\vskip 4pt
\begin{iband}[signal]
Akibat yang penting bagi pemakai: kemampuan itu \textbf{sekuat data pasangannya}. Jenis dokumen yang jarang muncul di internet --- formulir bank Indonesia, misalnya --- jauh lebih lemah daripada foto sehari-hari, dan itu bukan sesuatu yang bisa diperbaiki dengan prompt.
\end{iband}
\end{frame}

\framekicker{Penyandi}
\begin{frame}[shrink=25]{Petak, dan kenapa detail kecil hilang}
{\fontsize{9.4}{12.6}\selectfont\color{fg2}Gambar dipecah jadi petak berukuran tetap, dan tiap petak jadi satu vektor. Teks kecil yang jatuh di dalam satu petak bersama banyak hal lain harus berbagi satu vektor dengan semuanya.\par}\vskip 4pt
\begin{minipage}[t]{0.485\textwidth}
{\fontsize{9.4}{12.6}\selectfont\color{fg2}\textbf{Akibatnya}\par}\vskip 4pt
\begin{ibullets}
\item Angka kecil pada tabel padat sering salah baca
\item Menaikkan resolusi benar-benar menolong --- dan berlipat empat biayanya
\item Memotong bagian yang relevan sering lebih baik daripada menaikkan resolusi seluruh halaman
\end{ibullets}

\end{minipage}%
\hfill
\begin{minipage}[t]{0.485\textwidth}
{\fontsize{9.4}{12.6}\selectfont\color{fg2}\textbf{Yang bisa dilakukan}\par}\vskip 4pt
\begin{ibullets}
\item Deteksi wilayah dulu, lalu potong dan perbesar
\item Untuk tabel: alat khusus, bukan model umum
\item Untuk teks: OCR sungguhan mengalahkan model bahasa
\end{ibullets}

\end{minipage}%
\begin{iband}[signal]
Nasihat yang menghemat paling banyak waktu di alur dokumen: \hilite{pakai alat yang memang dibuat untuk itu}. Model bahasa yang melihat gambar bukan pengganti OCR.
\end{iband}
\end{frame}

\framekicker{Layar}
\begin{frame}[shrink=25]{Tangkapan layar: modalitas yang paling cepat tumbuh}
{\fontsize{9.4}{12.6}\selectfont\color{fg2}Agen yang mengoperasikan aplikasi lewat tangkapan layar menarik justru karena ia tidak butuh API --- ia memakai antarmuka yang sudah ada, sama seperti manusia.\par}\vskip 4pt
\begin{minipage}[t]{0.494\textwidth}
\begin{icard}
\cardh{Kenapa menggoda}
Sistem lama tanpa API tiba-tiba bisa diotomatiskan tanpa menunggu tim yang memeliharanya.
\end{icard}
\end{minipage}%
\hfill
\begin{minipage}[t]{0.494\textwidth}
\begin{icardbad}
\cardh{Kenapa rapuh}
Antarmuka berubah dan agennya rusak diam-diam. Tidak ada kontrak apa pun yang menjaganya.
\end{icardbad}
\end{minipage}%
{\fontsize{9.4}{12.6}\selectfont\color{fg2}Dan biayanya besar: tiap langkah butuh satu tangkapan layar baru, jadi \textbf{satu tugas sepuluh langkah berarti sepuluh gambar} di dalam satu percakapan yang sama.\par}\vskip 4pt
\begin{iband}[signal]
Urutan yang benar hampir selalu: \textbf{cari API-nya dulu, cari basis datanya, baru pertimbangkan layar.} Yang terakhir adalah pilihan ketika tidak ada pilihan lain.
\end{iband}
\end{frame}

\framekicker{Layar}
\begin{frame}[shrink=25]{Kalau memang harus lewat layar}
\begin{isteps}
\item \{'h': 'Batasi ke satu aplikasi, satu lingkungan', 'p': 'Agen yang bisa mengeklik apa saja di mesin mana saja adalah permukaan yang tidak bisa dijelaskan ke siapa pun.'\}
\item \{'h': 'Alat berupa aksi bernama, bukan ``klik di 300,450''', 'p': '\inlinecode{buka\_tagihan(id)} bisa diuji dan bertahan; koordinat piksel rusak pada perubahan tata letak pertama.'\}
\item \{'h': 'Tangkapan layar terakhir saja di konteks', 'p': 'Sepuluh tangkapan layar di riwayat adalah puluhan ribu token yang dibayar berulang.'\}
\item \{'h': 'Selalu di balik persetujuan untuk tindakan yang menulis', 'p': 'Sama seperti alat tulis mana pun. Klik yang mengirim uang tetap klik yang mengirim uang.'\}
\end{isteps}
\end{frame}

\coursesection{04}{Bentuk yang bekerja di dunia nyata}{Alat dulu, penalaran belakangan.}

\framekicker{Bentuk}
\begin{frame}[shrink=25]{Ekstrak dulu, nalarkan kemudian}
\begin{center}
\includegraphics[width=0.92\linewidth,height=0.60\textheight,keepaspectratio]{figs/agents09-pipe.tex.pdf}
\end{center}
\vskip 3pt{\color{ink3}\fontsize{7.4}{9.4}\selectfont Jalan yang lurus dan rapuh digambar sebagai garis putus.\par}
{\fontsize{9.4}{12.6}\selectfont\color{fg2}Bab 3 sudah menyebut kesimpulannya: penalaran berlangkah di atas gambar masih jauh lebih lemah daripada di atas teks. Jadi jangan minta satu panggilan mengerjakan keduanya.\par}\vskip 4pt
\begin{iband}[signal]
Bentuk yang bertahan: \textbf{alat mengubah gambar jadi data terstruktur; agen menalar atas data itu.} Hasilnya lebih murah, lebih tepat, dan --- yang paling penting --- bisa diperiksa.
\end{iband}
\end{frame}

\framekicker{Bentuk}
\begin{frame}[shrink=25]{Apa yang bisa dan tidak bisa diandalkan}
\begin{minipage}[t]{0.494\textwidth}
\begin{icardgood}
\cardh{Cukup andal}
Menjelaskan isi gambar, membaca teks tercetak yang jelas, mengenali jenis dokumen, menemukan bagian yang diminta.
\end{icardgood}
\end{minipage}%
\hfill
\begin{minipage}[t]{0.494\textwidth}
\begin{icardbad}
\cardh{Sering meleset}
Membaca angka pada tabel padat, mencocokkan nilai dengan kolomnya, tulisan tangan, stempel yang menimpa teks.
\end{icardbad}
\end{minipage}%


\vskip 4pt

\begin{minipage}[t]{0.494\textwidth}
\begin{icardbad}
\cardh{Jangan diandalkan}
Menghitung jumlah benda, membandingkan posisi, dan aritmetika atas angka yang dibacanya sendiri.
\end{icardbad}
\end{minipage}%
\hfill
\begin{minipage}[t]{0.494\textwidth}
\begin{icardgood}
\cardh{Serahkan ke kode}
Semua aritmetika. Model membaca angkanya; kode yang menghitungnya.
\end{icardgood}
\end{minipage}%
{\fontsize{9.4}{12.6}\selectfont\color{fg2}Kartu ketiga penting untuk dokumen keuangan: model yang membaca dua belas angka dengan benar lalu menjumlahkannya \textbf{tetap bisa salah pada jumlahnya} --- dan jumlah itu yang dipakai orang.\par}\vskip 4pt
\end{frame}

\framekicker{Bentuk}
\begin{frame}[shrink=25]{Ekstraksi yang bisa dipercaya menyebut keyakinannya}
\begin{isteps}
\item \{'h': 'Kembalikan nilai beserta lokasinya', 'p': 'Kotak koordinat tempat nilai itu terbaca. Ini yang membuat manusia bisa memeriksa dalam dua detik.'\}
\item \{'h': 'Bedakan kosong dari tidak terbaca', 'p': 'Medan yang memang kosong dan medan yang gagal dibaca menuntut tindakan yang berbeda sepenuhnya.'\}
\item \{'h': 'Sediakan ambang untuk minta bantuan', 'p': 'Di bawah ambang, eskalasikan --- jangan menebak. Bab 5: alat menyerah dengan alasan adalah hasil yang sah.'\}
\item \{'h': 'Simpan gambarnya, bukan di konteks', 'p': 'Di penyimpanan yang punya pemilik dan retensi, dengan pengenal yang bisa dirujuk jejak.'\}
\end{isteps}
\end{frame}

\framekicker{Dokumen}
\begin{frame}[shrink=25]{Tabel adalah kasus tersulit, dan yang paling sering dibutuhkan}
{\fontsize{9.4}{12.6}\selectfont\color{fg2}Membaca satu angka dari sebuah halaman relatif mudah. Membaca \textbf{tabel} menuntut dua hal sekaligus: mengenali nilainya, dan mengetahui nilai itu milik baris dan kolom yang mana.\par}\vskip 4pt
\begin{minipage}[t]{0.3253\textwidth}
\begin{icardbad}
\cardh{Kesalahan pergeseran}
Nilai benar, kolomnya salah. Ini yang paling berbahaya, sebab hasilnya masuk akal dan salah.
\end{icardbad}
\end{minipage}%
\hfill
\begin{minipage}[t]{0.3253\textwidth}
\begin{icard}
\cardh{Sel gabungan}
Judul yang membentang dua kolom, atau baris yang bersambung ke halaman berikutnya.
\end{icard}
\end{minipage}%
\hfill
\begin{minipage}[t]{0.3253\textwidth}
\begin{icard}
\cardh{Tabel lintas halaman}
Judul kolomnya ada di halaman sebelumnya, dan tiap halaman diproses terpisah.
\end{icard}
\end{minipage}%
\begin{iband}[signal]
Karena itu untuk tabel keuangan, \textbf{alat khusus pengurai tabel mengalahkan model umum} --- dan kalau tidak ada, minta model mengembalikan nilai \textbf{beserta nama baris dan kolomnya}, supaya pergeseran bisa ditangkap kode.
\end{iband}
\end{frame}

\framekicker{Dokumen}
\begin{frame}[shrink=25]{Tulisan tangan, tanda tangan, dan stempel}
\begin{minipage}[t]{0.485\textwidth}
{\fontsize{9.4}{12.6}\selectfont\color{fg2}\textbf{Tulisan tangan}\par}\vskip 4pt
{\fontsize{9.4}{12.6}\selectfont\color{fg2}Jauh lebih lemah daripada teks tercetak, dan lebih lemah lagi untuk angka. Untuk nilai yang menentukan keputusan, ini bukan jalur yang bisa dipercaya tanpa pemeriksaan manusia.\par}\vskip 4pt

\end{minipage}%
\hfill
\begin{minipage}[t]{0.485\textwidth}
{\fontsize{9.4}{12.6}\selectfont\color{fg2}\textbf{Tanda tangan dan stempel}\par}\vskip 4pt
{\fontsize{9.4}{12.6}\selectfont\color{fg2}Mendeteksi \textbf{keberadaannya} masuk akal dan berguna. Menilai \textbf{keasliannya} adalah tugas yang berbeda sepenuhnya, dan bukan tugas model bahasa.\par}\vskip 4pt

\end{minipage}%
\begin{iband}[signal]
Perbedaan di kolom kanan pantas ditulis di dokumen rancangan: ``ada tanda tangan di kotak itu'' adalah temuan yang berguna; \hilite{``tanda tangannya sah'' adalah klaim yang tidak boleh dibuat sistem ini}
\end{iband}
\end{frame}

\framekicker{Dokumen}
\begin{frame}[shrink=25]{Dokumen banyak halaman: satu keputusan yang menentukan biaya}
{\fontsize{9.4}{12.6}\selectfont\color{fg2}Berkas dua puluh halaman yang dikirim seluruhnya adalah puluhan ribu token --- dan biasanya sembilan belas halamannya tidak relevan.\par}\vskip 4pt
\begin{isteps}
\item \{'h': 'Klasifikasikan halaman dulu, dengan model kecil atau aturan', 'p': 'Halaman mana yang memuat neraca, mana yang lampiran. Murah, sekali, dan memotong sebagian besar biaya.'\}
\item \{'h': 'Proses halaman yang relevan saja, dengan resolusi penuh', 'p': 'Anggaran resolusi dibelanjakan di tempat yang benar.'\}
\item \{'h': 'Simpan pemetaannya', 'p': 'Nilai ini dari halaman 7 kotak sekian --- itu kutipan yang bisa diperiksa.'\}
\end{isteps}
\begin{iband}[signal]
Pola yang sama seperti pengambilan di Bab 4: \textbf{jaring lebar dengan yang murah, saringan halus dengan yang mahal.} Bentuk yang muncul lagi karena alasan yang sama.
\end{iband}
\end{frame}

\framekicker{Dokumen}
\begin{frame}[shrink=25]{Membandingkan dua pendekatan ekstraksi dengan angka}
\begin{itable}{>{\raggedright\arraybackslash}p{0.2444\textwidth}>{\raggedright\arraybackslash}p{0.3384\textwidth}>{\raggedright\arraybackslash}p{0.3572\textwidth}}
\thead{Pendekatan} & \thead{Kuat pada} & \thead{Lemah pada} \\ \midrule
OCR + aturan & Format tetap, volume besar, biaya nyaris nol & Variasi tata letak; tiap format baru butuh kerja \\
OCR + model bahasa & Tata letak bervariasi, medan yang bernama & Biaya per dokumen; tetap perlu validasi \\
Model penglihatan langsung & Dokumen yang tak terstruktur & Paling mahal, paling sulit diaudit, paling sering salah pada tabel \\
\textbf{Campuran} & Kenyataan di kebanyakan sistem & Perlu memutuskan mana yang dipakai kapan \\
\end{itable}
{\fontsize{9.4}{12.6}\selectfont\color{fg2}Keputusannya tidak perlu ditebak: \textbf{jalankan ketiganya pada dua puluh dokumen nyata} dan bandingkan ketepatan, biaya, dan laju eskalasi. Satu hari kerja, dan jawabannya tidak perlu diperdebatkan lagi.\par}\vskip 4pt
\end{frame}

\coursesection{05}{Bukti dan kewajiban}{Dua hal yang berubah paling besar, dan paling sering terlambat dipikirkan.}

\framekicker{Bukti}
\begin{frame}[shrink=25]{Jawaban dari gambar harus bisa ditunjuk}
\begin{center}
\includegraphics[width=0.92\linewidth,height=0.60\textheight,keepaspectratio]{figs/agents09-evidence.tex.pdf}
\end{center}
\vskip 3pt{\color{ink3}\fontsize{7.4}{9.4}\selectfont Satu pertanyaan yang memisahkan yang bisa diaudit dari yang tidak.\par}
{\fontsize{9.4}{12.6}\selectfont\color{fg2}Untuk teks, kutipan berarti pengenal potongan yang bisa diperiksa. Untuk gambar, bentuknya \textbf{lokasi} --- halaman, kotak, dan nilai yang terbaca di situ. Tanpa itu, satu-satunya cara memeriksa adalah membuka dokumennya dan mencari sendiri --- \hilite{agennya tidak menghemat pekerjaan siapa pun}\par}\vskip 4pt
\end{frame}

\framekicker{Kewajiban}
\begin{frame}[shrink=25]{Foto dokumen hampir selalu data pribadi}
{\fontsize{9.4}{12.6}\selectfont\color{fg2}Satu foto KTP atau rekening koran memuat lebih banyak data pribadi daripada seluruh permintaan teks yang menyertainya --- termasuk hal yang tidak diminta dan tidak dibutuhkan siapa pun.\par}\vskip 4pt
\begin{isteps}
\item \{'h': 'Gambar tidak keluar perimeter', 'p': 'Di demo, foto dari ponsel masuk ke penyimpanan bank dan berhenti di situ. \textbf{Tidak ada jalur kode} yang bisa mengirimnya ke penyedia model.'\}
\item \{'h': 'Kalau harus diproses model, batasi apa yang dikirim', 'p': 'Potongan yang relevan, bukan halaman penuh. Yang tidak dikirim tidak perlu dijelaskan.'\}
\item \{'h': 'Retensinya mengikuti dokumennya', 'p': 'Bukan mengikuti umur log. Ini yang paling sering salah.'\}
\end{isteps}
\begin{iband}[signal]
Pertanyaan yang akan diajukan pemeriksa, dan yang harus dijawab dengan \textbf{jalur kode}, bukan dengan kebijakan: \textit{ke mana perginya foto dokumen nasabah?}
\end{iband}
\end{frame}

\framekicker{Kewajiban}
\begin{frame}[shrink=25]{Gambar juga bisa membawa perintah}
{\fontsize{9.4}{12.6}\selectfont\color{fg2}Bab 4 menyebut perintah yang menumpang di dalam teks yang dibaca agen. Hal yang sama berlaku untuk gambar: teks di dalam gambar dibaca model, dan \textbf{model tidak membedakan} teks yang isinya data dari teks yang isinya perintah.\par}\vskip 4pt
\begin{minipage}[t]{0.494\textwidth}
\begin{icard}
\cardh{Bentuk seranganya}
Sebuah dokumen yang memuat kalimat ditujukan kepada agen, dicetak kecil atau ditempel di sudut halaman.
\end{icard}
\end{minipage}%
\hfill
\begin{minipage}[t]{0.494\textwidth}
\begin{icardaccent}
\cardh{Penanganannya sama}
Bukan di prompt. Di batas alat: apa yang boleh dipanggil sesudah membaca sesuatu yang tidak tepercaya.
\end{icardaccent}
\end{minipage}%
\begin{iband}[signal]
Dan satu tambahan khas gambar: kalau dokumennya berasal dari luar, \textbf{perlakukan hasil ekstraksinya sebagai masukan tidak tepercaya} --- divalidasi, bukan diteruskan apa adanya.
\end{iband}
\end{frame}

\coursesection{05b}{Modalitas lain}{Suara, video, dan pola yang sama muncul lagi.}

\framekicker{Modalitas}
\begin{frame}[shrink=25]{Satu ruang, banyak pintu masuk}
\begin{center}
\includegraphics[width=0.92\linewidth,height=0.60\textheight,keepaspectratio]{figs/agents09-modalities.tex.pdf}
\end{center}
\vskip 3pt{\color{ink3}\fontsize{7.4}{9.4}\selectfont Tiap modalitas punya penyandinya sendiri, semuanya bermuara ke satu ruang.\par}
{\fontsize{9.4}{12.6}\selectfont\color{fg2}Bentuknya seragam: tiap modalitas punya penyandi yang mengubahnya jadi vektor, lalu semuanya dipetakan ke ruang yang sama supaya model bahasa bisa memperlakukannya sebagai token.\par}\vskip 4pt
\begin{iband}[signal]
Karena bentuknya seragam, \textbf{pertanyaan rekayasanya juga seragam}: berapa token yang dihasilkan, apa yang hilang saat penyandian, dan apakah hasilnya bisa dikutip.
\end{iband}
\end{frame}

\framekicker{Suara}
\begin{frame}[shrink=25]{Suara: hampir selalu transkripsikan dulu}
{\fontsize{9.4}{12.6}\selectfont\color{fg2}Model bisa menerima suara langsung, dan untuk sebagian tugas itu berguna --- nada bicara, jeda, siapa yang berbicara. Untuk hampir semua agen bisnis, jalur yang benar tetap transkripsi lewat alat.\par}\vskip 4pt
\begin{minipage}[t]{0.3253\textwidth}
\begin{icardgood}
\cardh{Bisa disimpan dan dicari}
Transkrip masuk basis data, bisa dicari, bisa dikutip. Gelombang suara tidak.
\end{icardgood}
\end{minipage}%
\hfill
\begin{minipage}[t]{0.3253\textwidth}
\begin{icardgood}
\cardh{Jauh lebih murah}
Satu jam rapat sebagai teks adalah beberapa ribu token; sebagai suara jauh lebih besar.
\end{icardgood}
\end{minipage}%
\hfill
\begin{minipage}[t]{0.3253\textwidth}
\begin{icard}
\cardh{Yang hilang}
Nada, keraguan, siapa berbicara. Kalau itu penting, catat sebagai metadata --- jangan bawa gelombangnya.
\end{icard}
\end{minipage}%
\begin{iband}[signal]
Dan satu kewajiban yang sering terlewat: \textbf{rekaman suara orang adalah data pribadi}, dan di banyak konteks perlu izin yang berbeda dari izin memproses teksnya.
\end{iband}
\end{frame}

\framekicker{Video}
\begin{frame}[shrink=25]{Video: biaya berlipat dengan durasi, dan hampir selalu bisa dihindari}
{\fontsize{9.4}{12.6}\selectfont\color{fg2}Video disandikan sebagai rentetan bingkai, kadang dengan tambahan penanda waktu. Artinya biayanya adalah biaya gambar \textbf{dikalikan jumlah bingkai} --- dan satu menit pada satu bingkai per detik sudah enam puluh gambar.\par}\vskip 4pt
\begin{isteps}
\item \{'h': 'Ambil bingkai kunci, bukan semuanya', 'p': 'Saat berubah, saat ada gerakan, atau satu per beberapa detik. Sering satu bingkai per adegan sudah cukup.'\}
\item \{'h': 'Pisahkan suaranya dan transkripsikan', 'p': 'Untuk sebagian besar video, transkrip menjawab lebih banyak pertanyaan daripada bingkainya, dengan biaya jauh lebih kecil.'\}
\item \{'h': 'Simpan penanda waktu', 'p': 'Supaya jawabannya bisa menunjuk detik keberapa --- kutipan, dalam bentuk waktu.'\}
\end{isteps}
\end{frame}

\framekicker{Modalitas}
\begin{frame}[shrink=25]{Apa yang hilang di tiap penyandian}
\begin{itable}{>{\raggedright\arraybackslash}p{0.2068\textwidth}>{\raggedright\arraybackslash}p{0.3572\textwidth}>{\raggedright\arraybackslash}p{0.376\textwidth}}
\thead{Modalitas} & \thead{Yang bertahan} & \thead{Yang hilang} \\ \midrule
Gambar & Isi, tata letak kasar, teks besar & Angka kecil, hubungan baris-kolom, detail di bawah ukuran petak \\
Suara & Kata-kata & Nada, keraguan, siapa berbicara --- kecuali dicatat terpisah \\
Video & Isi bingkai yang diambil & Semua yang terjadi di antara bingkai yang diambil \\
Tabel sebagai gambar & Angkanya & \textbf{Hubungan angka dengan kolomnya} --- dan itu yang dipakai \\
\end{itable}
\begin{iband}[signal]
Baris terakhir yang paling sering menyebabkan kesalahan yang tidak terlihat: nilainya benar, kolomnya salah, dan hasilnya tetap terlihat masuk akal.
\end{iband}
\end{frame}

\framekicker{Kewajiban}
\begin{frame}[shrink=25]{Apa yang boleh disimpan dari sebuah dokumen}
{\fontsize{9.4}{12.6}\selectfont\color{fg2}Godaan yang wajar sesudah ekstraksi berhasil: simpan semuanya, siapa tahu berguna. Itu cara paling cepat menciptakan basis data data pribadi yang tidak ada di rancangan siapa pun.\par}\vskip 4pt
\begin{itable}{>{\raggedright\arraybackslash}p{0.282\textwidth}>{\raggedright\arraybackslash}p{0.2068\textwidth}>{\raggedright\arraybackslash}p{0.4512\textwidth}}
\thead{Yang diekstrak} & \thead{Simpan?} & \thead{Alasannya} \\ \midrule
Angka yang dipakai keputusan & Ya & Beserta lokasinya --- ini bukti yang dibutuhkan pemeriksa \\
Jenis dan tanggal dokumen & Ya & Metadata, dan diperlukan untuk retensi \\
Nama, alamat, nomor identitas & \textbf{Tidak} & Agen tidak membutuhkannya; petugas melihatnya di sistem inti \\
Gambar aslinya & Di penyimpanan dokumen & Dengan pemilik dan retensi dokumen, bukan retensi log \\
\end{itable}
\begin{iband}[signal]
Baris ketiga adalah keputusan rancangan, bukan pembatasan yang merepotkan: \textbf{yang tidak disimpan tidak perlu dijaga, tidak perlu dihapus, dan tidak bisa bocor.}
\end{iband}
\end{frame}

\coursesection{06}{Praktik}{Merancang, mengukur, dan menguji alur yang menerima gambar.}

\framekicker{Praktik}
\begin{frame}[shrink=25]{Urutan merancang alur dokumen}
\begin{isteps}
\item \{'h': 'Cari tahu apakah teksnya sudah ada', 'p': 'PDF digital, hasil sistem lain, berkas asal. Kalau ada, seluruh bab ini bisa dilewati.'\}
\item \{'h': 'Tulis alat ekstraksi, bukan prompt ekstraksi', 'p': 'Keluarannya berskema, dengan lokasi dan tingkat keyakinan.'\}
\item \{'h': 'Nalarkan atas data terstrukturnya', 'p': 'Di situ semua yang dipelajari tujuh bab sebelumnya berlaku seperti biasa.'\}
\item \{'h': 'Simpan gambarnya di luar konteks', 'p': 'Dengan pengenal, pemilik, dan retensi.'\}
\end{isteps}
\begin{iband}[signal]
Langkah pertama yang paling sering dilewat, dan paling sering menghapus seluruh masalah: \textbf{banyak dokumen yang difoto sebenarnya tersedia dalam bentuk teks di sistem lain.}
\end{iband}
\end{frame}

\framekicker{Praktik}
\begin{frame}[shrink=25]{Mengujinya: sebagian besar tidak butuh model}
\begin{minipage}[t]{0.494\textwidth}
\begin{icardgood}
\cardh{Uji skema ekstraksi}
Medan yang wajib ada, tipe yang benar, lokasi yang masuk akal. Kode biasa.
\end{icardgood}
\end{minipage}%
\hfill
\begin{minipage}[t]{0.494\textwidth}
\begin{icardgood}
\cardh{Kumpulan gambar tetap}
Dua puluh dokumen nyata dengan nilai yang benar sudah dicatat. Ini kumpulan uji Bab 7, dalam bentuk gambar.
\end{icardgood}
\end{minipage}%


\vskip 4pt

\begin{minipage}[t]{0.494\textwidth}
\begin{icard}
\cardh{Kasus yang buruk sengaja}
Miring, gelap, terpotong, stempel menimpa. Di sinilah sistem menunjukkan apakah ia tahu ia tidak tahu.
\end{icard}
\end{minipage}%
\hfill
\begin{minipage}[t]{0.494\textwidth}
\begin{icardgood}
\cardh{Uji jalur data}
Pastikan tidak ada jalur kode dari penyimpanan gambar ke penyedia model. Uji struktural, seperti Bab 5.
\end{icardgood}
\end{minipage}%
\end{frame}

\framekicker{Praktik}
\begin{frame}[shrink=25]{Angka yang pantas dipantau pada alur bergambar}
\begin{itable}{>{\raggedright\arraybackslash}p{0.2632\textwidth}>{\raggedright\arraybackslash}p{0.2632\textwidth}>{\raggedright\arraybackslash}p{0.4136\textwidth}}
\thead{Angka} & \thead{Sehat kalau} & \thead{Gejala kalau tidak} \\ \midrule
Token gambar per tugas & Datar dan kecil & Naik $\rightarrow$ gambar tidak dikeluarkan dari riwayat \\
Laju ekstraksi gagal & Stabil & Naik $\rightarrow$ mutu foto turun, atau jenis dokumen berubah \\
Laju eskalasi karena tak terbaca & Ada, dan kecil & Nol $\rightarrow$ ambang keyakinannya tidak berfungsi \\
Selisih nilai terbaca lawan sistem & Mendekati nol & Ada $\rightarrow$ ekstraksinya salah, atau dokumennya bukan yang dikira \\
\end{itable}
{\fontsize{9.4}{12.6}\selectfont\color{fg2}Baris ketiga sekali lagi: \textbf{nol eskalasi bukan keberhasilan.} Dokumen yang buruk pasti ada; sistem yang tidak pernah mengeluh sedang menebak.\par}\vskip 4pt
\end{frame}

\framekicker{Praktik}
\begin{frame}[shrink=25]{Waktu tunggu pada alur bergambar}
{\fontsize{9.4}{12.6}\selectfont\color{fg2}Gambar menambah waktu di dua tempat sekaligus: mengunggah berkasnya, dan memproses masukan yang jauh lebih besar. Pada ponsel di jaringan lapangan, yang pertama sering lebih besar daripada yang kedua.\par}\vskip 4pt
\begin{minipage}[t]{0.494\textwidth}
\begin{icardgood}
\cardh{Perkecil sebelum dikirim}
Di sisi ponsel, bukan di peladen. Resolusi yang cukup untuk dibaca sering jauh di bawah resolusi kamera.
\end{icardgood}
\end{minipage}%
\hfill
\begin{minipage}[t]{0.494\textwidth}
\begin{icard}
\cardh{Kerjakan yang lain sambil menunggu}
Ekstraksi dokumen tidak menghalangi pengambilan data lain --- itu langkah yang boleh sejajar.
\end{icard}
\end{minipage}%
\begin{iband}[signal]
Dan satu hal yang lebih besar pengaruhnya daripada keduanya: \textbf{tunjukkan apa yang sedang terjadi.} Layar yang menyebut ``membaca halaman 3 dari 5'' mengubah pengalaman menunggu lebih banyak daripada penghematan satu detik.
\end{iband}
\end{frame}

\framekicker{Praktik}
\begin{frame}[shrink=25]{Enam kegagalan khas alur bergambar}
\begin{minipage}[t]{0.3253\textwidth}
\begin{icardbad}
\cardh{Gambar menetap di riwayat}
Dibayar ulang tiap giliran. Penyebab pembengkakan biaya nomor satu di alur dokumen.
\end{icardbad}
\end{minipage}%
\hfill
\begin{minipage}[t]{0.3253\textwidth}
\begin{icardbad}
\cardh{Model menghitung sendiri}
Membaca dua belas angka dengan benar lalu salah menjumlahkannya.
\end{icardbad}
\end{minipage}%
\hfill
\begin{minipage}[t]{0.3253\textwidth}
\begin{icardbad}
\cardh{Nilai bergeser kolom}
Masuk akal, dan salah. Tidak tertangkap tanpa nama baris dan kolom.
\end{icardbad}
\end{minipage}%


\vskip 4pt

\begin{minipage}[t]{0.3253\textwidth}
\begin{icardbad}
\cardh{Tak terbaca dianggap kosong}
Medan gagal dibaca diperlakukan seperti medan yang memang kosong.
\end{icardbad}
\end{minipage}%
\hfill
\begin{minipage}[t]{0.3253\textwidth}
\begin{icardbad}
\cardh{Gambar keluar perimeter}
Dikirim ke penyedia model karena tidak ada yang melarangnya di kode.
\end{icardbad}
\end{minipage}%
\hfill
\begin{minipage}[t]{0.3253\textwidth}
\begin{icardbad}
\cardh{Resolusi penuh untuk semuanya}
Empat kali biaya untuk ketepatan yang sering tidak berubah.
\end{icardbad}
\end{minipage}%
\end{frame}

\framekicker{Praktik}
\begin{frame}[shrink=25]{Kapan multi-modal adalah alat yang salah}
\begin{minipage}[t]{0.485\textwidth}
{\fontsize{9.4}{12.6}\selectfont\color{fg2}\textbf{Pakai alat lain}\par}\vskip 4pt
\begin{ibullets}
\item Teksnya sudah ada di sistem lain
\item Formatnya tetap dan volumenya besar $\rightarrow$ OCR + aturan
\item Yang dibutuhkan hitungan $\rightarrow$ kode, selalu
\item Keasliannya yang dipertanyakan $\rightarrow$ bukan tugas model bahasa
\end{ibullets}

\end{minipage}%
\hfill
\begin{minipage}[t]{0.485\textwidth}
{\fontsize{9.4}{12.6}\selectfont\color{fg2}\textbf{Di sinilah ia menang}\par}\vskip 4pt
\begin{ibullets}
\item Tata letak bervariasi dan medannya bernama
\item Dokumen yang tidak terstruktur
\item Menentukan jenis dokumen sebelum diproses
\item Menjelaskan isi gambar kepada manusia
\end{ibullets}

\end{minipage}%
\begin{iband}[signal]
Kolom kiri baris pertama tetap yang paling sering benar dan paling sering terlewat: \textbf{tanyakan dulu apakah dokumen itu pernah ada dalam bentuk teks} sebelum membangun apa pun.
\end{iband}
\end{frame}

\framekicker{Praktik}
\begin{frame}[shrink=25]{Yang dipakai demo, dan apa yang sengaja tidak ada di sana}
{\fontsize{9.4}{12.6}\selectfont\color{fg2}Aplikasi lapangan di demo memotret dokumen pengajuan. Yang menarik dari rancangannya bukan apa yang dilakukannya, melainkan \textbf{apa yang tidak mungkin dilakukannya.}\par}\vskip 4pt
\begin{isteps}
\item \{'h': 'Foto pergi ke penyimpanan bank, dan berhenti di situ', 'p': 'Tidak ada jalur kode dari penyimpanan itu ke penyedia model mana pun.'\}
\item \{'h': 'Agen menerima pengenal pengajuan, bukan gambar', 'p': 'Alat yang dimilikinya tidak mengembalikan gambar, jadi ia tidak bisa meneruskannya ke mana pun.'\}
\item \{'h': 'Klaimnya diuji, bukan dijanjikan', 'p': '\inlinecode{test\_the\_agents\_tools\_cannot\_reach\_personal\_data} gagal kalau seseorang menambahkan alat yang melanggarnya.'\}
\end{isteps}
\begin{iband}[signal]
Ini penerapan langsung prinsip Bab 1 pada modalitas baru: \textbf{batas kemampuan ada di ketiadaan alat}, dan ketiadaan itu bisa diuji.
\end{iband}
\end{frame}

\framekicker{Praktik}
\begin{frame}[shrink=25]{Urutan memasangnya, dari yang paling murah}
\begin{isteps}
\item \{'h': 'Cari teksnya', 'p': 'Kalau ada, selesai. Ini langkah yang paling sering menghapus seluruh pekerjaan.'\}
\item \{'h': 'OCR sungguhan untuk teks tercetak', 'p': 'Lebih murah dan lebih tepat daripada model bahasa untuk pekerjaan ini.'\}
\item \{'h': 'Model penglihatan untuk yang tata letaknya bervariasi', 'p': 'Dengan keluaran berskema, lokasi, dan tingkat keyakinan.'\}
\item \{'h': 'Manusia untuk yang di bawah ambang', 'p': 'Dan hitung berapa sering itu terjadi --- angka itu memberi tahu apakah alurnya layak dipakai.'\}
\end{isteps}
\begin{iband}[signal]
Perhatikan bahwa dua langkah pertama tidak melibatkan model bahasa sama sekali, dan keduanya menyelesaikan sebagian besar dokumen yang benar-benar dihadapi bank.
\end{iband}
\end{frame}

\framekicker{Penutup}
\begin{frame}[shrink=25]{Yang dibawa pulang dari bab ini}
\begin{isteps}
\item \{'h': 'Gelungnya sama; biaya, bukti, dan kewajibannya berubah', 'p': 'Tidak ada mekanisme baru --- hanya tiga hal yang jadi jauh lebih besar.'\}
\item \{'h': 'Satu gambar bukan satu lampiran', 'p': 'Foto dokumen $\approx$ 3,3$\times$ halaman teks yang sama, dan biayanya tumbuh menurut luas.'\}
\item \{'h': 'Ekstrak dengan alat, nalarkan atas datanya', 'p': 'Lebih murah, lebih tepat, dan satu-satunya bentuk yang bisa diperiksa.'\}
\item \{'h': 'Jawaban dari gambar harus bisa ditunjuk', 'p': 'Halaman, kotak, nilai. Tanpa itu tidak ada yang bisa diaudit.'\}
\item \{'h': 'Foto dokumen adalah data pribadi', 'p': 'Dan pertanyaan \textbackslash\{\}\textbackslash\{\}u201cke mana perginya\textbackslash\{\}\textbackslash\{\}u201d harus dijawab dengan jalur kode.'\}
\end{isteps}
\end{frame}
\note{Pertanyaan penutup: dokumen yang kalian foto, apakah teksnya sebenarnya sudah ada di suatu sistem? Jawabannya cukup sering ya, dan itu menghapus seluruh bab ini.}

\end{document}
